% Math macros for the Topping blueprint — consolidated for hgraph/KaTeX.

% --- from macros/common.tex ---
\newtheorem{theorem}{Theorem}[chapter]
\newtheorem{proposition}[theorem]{Proposition}
\newtheorem{lemma}[theorem]{Lemma}
\newtheorem{corollary}[theorem]{Corollary}
\newtheorem{conjecture}[theorem]{Conjecture}

\theoremstyle{definition}
\newtheorem{definition}[theorem]{Definition}
\newtheorem{remark}[theorem]{Remark}
\newtheorem{example}[theorem]{Example}
\newtheorem{notation}[theorem]{Notation}
\newtheorem{convention}[theorem]{Convention}

\newcommand{\Ric}{\operatorname{Ric}}
\newcommand{\Rm}{\operatorname{Rm}}
\newcommand{\inj}{\operatorname{inj}}
\newcommand{\Vol}{\operatorname{Vol}}
\newcommand{\Hess}{\operatorname{Hess}}
\newcommand{\Isom}{\operatorname{Isom}}
\newcommand{\Lie}{\mathcal{L}}
\newcommand{\dG}{\delta G}
\newcommand{\Int}{\operatorname{int}}
\newcommand{\dd}{\,\mathrm{d}}
\newcommand{\M}{\mathcal{M}}
\newcommand{\id}{\operatorname{id}}
\newcommand{\tr}{\operatorname{tr}}
\newcommand{\Sym}{\operatorname{Sym}}
\newcommand{\RR}{\mathbb{R}}
\newcommand{\NN}{\mathbb{N}}
\newcommand{\ZZ}{\mathbb{Z}}
\newcommand{\e}{\varepsilon}
\newcommand{\abs}[1]{\lvert #1\rvert}

% Horizon metadata is parsed into the DAG and hidden in print output.
\providecommand{\dcref}[1]{}
\providecommand{\group}[1]{}
\providecommand{\source}[1]{}
\providecommand{\home}[1]{}
\providecommand{\github}[1]{}
\providecommand{\dochome}[1]{}
\providecommand{\uses}[1]{}
